\section*{الفصل \ref{ch:b2:plant-flowering} --- النمو التكاثري والتحكم في الإزهار}

\begin{solution}{exo:b2:plant-flowering:1}
\omterm{def:b2:plant-meristems:meristem}{المرستيم} الخضري: يصنع أوراقًا ببراعم إبطية، بلا نهاية
(فهو غير محدد). \omterm{def:b2:plant-flowering:transition}{والمرستيم النوري}: يصنع قنابات ومرستيمات
زهرية في آباطها، وهو غير محدد في أنواع كثيرة. \omterm{def:b2:plant-meristems:meristem}{والمرستيم}
الزهري: يصنع الدوّارات الأربع ويتوقف --- فهو محدد.
\end{solution}

\begin{solution}{exo:b2:plant-flowering:2}
النباتات القصيرة النهار تزهر حين يكون النهار أقصر من طول
حرج (الأقحوان، وفول الصويا، والأرز)؛ والطويلة النهار حين يكون
أطول (السبانخ، والقمح، و\emph{Arabidopsis})؛ والمحايدة النهار
تزهر حين تبلغ قياسًا (البندورة، والذرة). \omterm{prop:b2:plant-flowering:photoperiod}{والنبات القصير النهار}
يقيس طول الليل غير المنقطع.
\end{solution}

\begin{solution}{exo:b2:plant-flowering:3}
النمط البري: سبلة، وبتلة، \omterm{def:b2:angiosperm-reproduction:flower}{وسداة}، \omterm{def:b2:angiosperm-reproduction:flower}{وكربلة}. وبلا A: \omterm{def:b2:angiosperm-reproduction:flower}{كربلة}، \omterm{def:b2:angiosperm-reproduction:flower}{وسداة}، و
\omterm{def:b2:angiosperm-reproduction:flower}{سداة}، \omterm{def:b2:angiosperm-reproduction:flower}{وكربلة}. وبلا B: سبلة، وسبلة، \omterm{def:b2:angiosperm-reproduction:flower}{وكربلة}، \omterm{def:b2:angiosperm-reproduction:flower}{وكربلة}. وبلا C: سبلة، و
بتلة، وبتلة، وسبلة، وتتكرر \omterm{def:b2:angiosperm-reproduction:flower}{الزهرة} داخلها.
\end{solution}

\begin{solution}{exo:b2:plant-flowering:4}
أسابيع من البرد والندى (فالشتاء مضى)؛ والضوء (\omterm{def:b2:angiosperm-reproduction:seed}{فالبذرة} تقع
قرب السطح)؛ وتناوب الحرارات (فالتربة عارية بلا مظلة)؛
وغسل المطر (فالأمطار جاءت)؛ والخدش (بالمرور
في أمعاء أو بالاحتكاك في التربة)؛ والدخان والحرارة (فحريق
نظّف الأرض).
\end{solution}

\begin{solution}{exo:b2:plant-flowering:5}
16/8: ليل من 8 ساعات، دون الحرج 10: فلا. و12/12: نعم.
و12 ضوءًا و6 ظلامًا وومضة و6 ظلامًا: أي ليلتان من 6 ساعات: فلا.
و8 ضوءًا و4 ظلامًا و4 ضوءًا و8 ظلامًا: أطول ليل 8 ساعات: فلا.
\end{solution}

\begin{solution}{exo:b2:plant-flowering:6}
\omterm{def:b2:plant-flowering:florigen}{الفلوريجين} إشارة تنتقل بالتطعيم تُصنع في الورقة وهي
نفسها في الأنواع القصيرة النهار والطويلة النهار --- والفرق في متى
تصنعها الورقة، لا في ماذا تصنع. ومع قطع اللحاء
لا يستحث الطُّعم الإزهار: فالإشارة تسافر في
اللحاء.
\end{solution}

\begin{solution}{exo:b2:plant-flowering:7}
$n = \ln(1/0.15)/0.05 = 1.90/0.05 = 38$ يومًا باردًا. ومجموع الموجات
$47$: فيُرباع النبات في اليوم الحادي عشر من الموجة الثالثة.
ومع $\alpha = 0.02$ كان ليحتاج 95 يومًا ولما رُبّع.
\end{solution}

\begin{solution}{exo:b2:plant-flowering:8}
مورّثة الكابت تُسكت بعلامات صبغين تُنسخ
عند كل تقسم خيطي، فتتذكر كل خلية من الفرع الشتاء؛
أما في السلالة الجنسية والجنين الباكر فتُمحى العلامات وتُعبَّر المورّثة
من جديد. وهذا مفيد لأن \omterm{def:b2:angiosperm-reproduction:seed}{بذرة} تُنثر في الصيف يجب ألا
ترث إذن أبيها بالإزهار: بل عليها أن تنتظر
شتاءها هي.
\end{solution}

\begin{solution}{exo:b2:plant-flowering:9}
أحمر: تنبت (إذ يتحول الفيتوكروم إلى الصورة الماصة للأحمر البعيد،
Pfr، وهي الفعالة). وأحمر ثم أحمر بعيد: لا (فيعود إلى Pr).
وأحمر ثم أحمر بعيد ثم أحمر: تنبت. وأحمر ثم أحمر بعيد ثم أحمر ثم أحمر بعيد: لا. فالومضة الأخيرة وحدها هي التي تُحسب،
لأن كل تحويل تام \omterm{def:b2:angiosperm-reproduction:seed}{والبذرة} تقرأ الحالة النهائية
للصباغ.
\end{solution}

\begin{solution}{exo:b2:plant-flowering:10}
$p = 0.7$: $r = 1.28$ و$1.42$ و$1.40$ و$-\infty$ من أجل $g = 0.4, 0.6, 0.8,
1$: فالأفضل قرب $g = 0.7$. و$p = 0.95$: $1.99$ و$2.33$ و$2.55$
و$-\infty$: فينتقل الأمثل نحو $g = 1$ (نحو $0.948$)، أي احتياطي
رمزي فقط. فالصحارى، حيث تندر السنوات الجيدة، تفضّل سكونًا ثقيلًا
وبنك بذور كبيرًا؛ والمروج، حيث معظم السنوات جيد،
تفضّل إنبات كل شيء تقريبًا.
\end{solution}

\begin{solution}{exo:b2:plant-flowering:11}
B في كل الدوّارات: بتلة، وبتلة، \omterm{def:b2:angiosperm-reproduction:flower}{وسداة}، \omterm{def:b2:angiosperm-reproduction:flower}{وسداة}. وC في كل الدوّارات (إذ يكبت C المورّثة A): \omterm{def:b2:angiosperm-reproduction:flower}{كربلة}، \omterm{def:b2:angiosperm-reproduction:flower}{وسداة}، \omterm{def:b2:angiosperm-reproduction:flower}{وسداة}، \omterm{def:b2:angiosperm-reproduction:flower}{وكربلة}. وبلا A وبلا B: يكون C وحده
في كل مكان --- أي كرابل في الدوّارات الأربع. وC ينهي
\omterm{def:b2:plant-meristems:meristem}{المرستيم} أيضًا؛ ودونه يبقى المركز غير محدد ويصنع \omterm{def:b2:angiosperm-reproduction:flower}{زهرة}
داخل \omterm{def:b2:angiosperm-reproduction:flower}{زهرة}.
\end{solution}

\begin{solution}{exo:b2:plant-flowering:12}
طول النهار يُكامَل عبر الليالي بتزامن الضوء
مع الساعة، ويحتاج الاستحثاث إلى عدة ليال كهذه؛ والبرد
يُكامَل علاماتِ صبغين متراكمة؛ وبنك البذور
يكامل عبر السنوات بتوزيع الإنبات. والمكامل
يتجاهل يومًا شاذًا واحدًا، أو ذوبان كانون الثاني، أو سنة رديئة واحدة؛
أما العتبة على القيمة اللحظية فكانت لتُزهر النبات في
أسبوع دافئ من شباط، أو تُنبت كل \omterm{def:b2:angiosperm-reproduction:seed}{بذرة} في أول مطرة.
\end{solution}

\begin{solution}{pb:b2:plant-flowering:1}
\textbf{1.} ليل من 8 ساعات أقصر من الحرج 9.5 ساعة:
فلا إزهار. وليبيع في آب على المزارع أن يطيل الليالي
اصطناعيًا بالقماش الأسود.
\textbf{2.} ليل من 12 ساعة. والتفتح في 15 آب يعني استحثاثًا
اكتمل قبل 8 أسابيع، أي نحو 20 حزيران؛
فيجب أن تبدأ الليالي الاستحثاثية الاثنتا عشرة نحو 8 حزيران.
\textbf{3.} أضئ المحصول بضع دقائق في منتصف كل
ليلة: فيصير ليل 14.5 ساعة ليلين من نحو 7 ساعات، وكلاهما دون
الطول الحرج. وتكفي الومضة لأن الفيتوكروم
يتحول في دقائق والنبات لا يقيس إلا طول
الظلام غير المنقطع.
\textbf{4.} الومضة عند 21:00 تترك فترة ظلام من
21:00 إلى 07:30، أي 10.5 ساعة --- وهي أطول من 9.5:
فيستحث ذلك المقطع. ويجب أن يقع القطع بحيث لا يتجاوز أي مقطع
الطول الحرج.
\textbf{5.} ليلة واحدة فائتة تكسر سلسلة الليالي الاثنتي عشرة
الاستحثاثية المتتالية: فيعود العدّ إلى البداية ويتأخر الإزهار
نحو ما تراكم من أيام. أما الخرشوف الشوكي فيحتاج إلى
ليلة استحثاثية واحدة وكان ليكون قد التزم سلفًا.
\textbf{6.} الفيتوكروم لا يمتص الضوء الأخضر، فلا يرى النبات
الومضة؛ والأحمر يحوله إلى الصورة الفعالة التي
تُقرأ نهارًا؛ والأحمر البعيد بعده مباشرة يعيده، فلا
ينقطع الليل.
\textbf{7.} يزهر \omterm{prop:b2:plant-flowering:photoperiod}{النبات الطويل النهار}: \omterm{def:b2:plant-flowering:florigen}{فالفلوريجين} يعبر وصلة
التطعيم في اللحاء. وهذا يعيّن إشارة متنقلة تنتقل بالتطعيم
مشتركة بين صنفي طول النهار.
\textbf{8.} $n = \ln 10/0.06 = 38.4$: أي 39 يومًا باردًا.
\textbf{9.} 10 في تشرين الثاني، و35 بنهاية كانون الثاني، ويقع اليوم البارد
39 في اليوم البارد الرابع من شباط: فيُرباع في
مطلع شهر شباط.
\textbf{10.} ستة أيام: $f = e^{-0.36} = 0.70 > 0.10$: فيبقى
خضريًا طوال الصيف ولا يعطي حبًا --- فقمح الشتاء يجب أن
يُزرع في الخريف.
\textbf{11.} دون الكابت يزهر بلا برد: فيمكن
زرعه في الربيع وحصاده في الصيف نفسه --- أي قمح ربيعي.
\textbf{12.} علامات الإسكات على صبغين الكابت
تُنسخ عند كل تقسم خيطي، فتتذكر القمة عبر الربيع؛ وفي
\omterm{def:b2:meiosis-heredity:meiosis}{التقسم المنصّف} والجنين تُمحى العلامات، وعلى كل جيل
أن يعدّ شتاءه هو.
\textbf{13.} يجب أن تقع الذاكرة في الخلايا التي ستصنع
\omterm{def:b2:angiosperm-reproduction:flower}{الزهرة} والتي تبقى عبر الشتاء --- أي \omterm{def:b2:plant-meristems:meristem}{المرستيم} --- وهي
حالة من صبغينها لا إشارة متنقلة؛ أما الأوراق التي
تقيس طول النهار فتُسقط، وتصدّر بروتينًا بدل ذلك.
\textbf{14.} الصنف C: سبلة، وبتلة، وبتلة، وسبلة (فينتشر A إلى
الدوّارة الرابعة وتتكرر \omterm{def:b2:angiosperm-reproduction:flower}{الزهرة} إلى الداخل).
\textbf{15.} لا سدى، فلا \omterm{def:b2:plant-life-cycles:heterospory}{لقاح}؛ وتتكون بضع كرابل حيث يكون C
فعالًا ضعيفًا، فتعقد بذور عارضة.
\textbf{16.} الصنف A: \omterm{def:b2:angiosperm-reproduction:flower}{كربلة}، \omterm{def:b2:angiosperm-reproduction:flower}{وسداة}، \omterm{def:b2:angiosperm-reproduction:flower}{وسداة}، \omterm{def:b2:angiosperm-reproduction:flower}{وكربلة}.
\textbf{17.} C يكبت A: فتصير الدوّارة 1 \omterm{def:b2:angiosperm-reproduction:flower}{كربلة}، والدوّارة 2 (حيث B وC)
\omterm{def:b2:angiosperm-reproduction:flower}{سداة} --- أي \omterm{def:b2:angiosperm-reproduction:flower}{كربلة}، \omterm{def:b2:angiosperm-reproduction:flower}{وسداة}، \omterm{def:b2:angiosperm-reproduction:flower}{وسداة}، \omterm{def:b2:angiosperm-reproduction:flower}{وكربلة}، وهي \omterm{def:b2:angiosperm-reproduction:flower}{زهرة} الطافرة A.
\textbf{18.} مورّثات ABC أسرة واحدة من \omterm{prop:b2:cell-differentiation:transcriptional}{عوامل الاستنساخ}
موروثة من السلف المشترك لكل النباتات الزهرية؛
فالوحدة النمطية كانت في مكانها قبل أن يتباعد الورد وحنك السبع و\emph{Arabidopsis}،
وقد أبقاها كل منها.
\textbf{19.} B دون A أو C لا يحدد عضوًا زهريًا: بل عضوًا شبيهًا
بالورقة أو بالقنابة.
\textbf{20.} $p = 0.8$: $r = 1.44$ (عند $g = 0.5$)، و$1.59$ (عند $0.7$)، و$1.55$
(عند $0.9$).
\textbf{21.} من الثلاثة، $g = 0.7$ (والأمثل الحقيقي قرب
$0.8$)؛ وعند $g = 1$ يكون عامل السنة الرديئة صفرًا وتُفنى السلالة
بأول سنة رديئة.
\textbf{22.} $p = 0.5$: $0.51$ و$0.41$ و$-0.03$: والقيمة $g = 0.5$ أفضلها،
وأقل منها كان ليكون أفضل --- أي \omterm{def:b2:angiosperm-reproduction:seed}{سكون} أكثر في الصحراء.
\textbf{23.} بذور المزارع مخزونة لا مدفونة، وتنجو
يقينًا؛ فلا يحتاج إلى تحوّط داخل الحقل، لكن عليه أن يبقي
احتياطيًا في المخزن ليعيد الزرع بعد سنة فاشلة --- وهو
الاحتياطي نفسه، يحفظه هو بدل أن تحفظه التربة.
\textbf{24.} \omterm{def:b2:angiosperm-reproduction:seed}{البذرة} الصغيرة تحمل مؤونة لفرع بطول
سنتيمتر أو سنتيمترين؛ ولو أنبتت في الظلام تحت التربة لماتت
قبل أن تبلغ الضوء، فهي تنتظر الضوء --- الذي لا يبلغها
إلا قرب السطح. والحرث يرفع البذور المدفونة إلى الأعلى و
يطلق دفقة من الأعشاب.
\textbf{25.} الليل الحرج 9.5 ساعة، والتغطية من نحو 8 حزيران
للتفتح في 15 آب؛ و39 يومًا باردًا؛ وأفضل $g \approx 0.8$ من أجل $p = 0.8$
ونحو $0.5$ من أجل $p = 0.5$.
\end{solution}
